\documentclass{article}
\usepackage[utf8]{inputenc}
\usepackage[T1]{fontenc}
\usepackage{geometry}
\geometry{
    a4paper,
    left=2.5cm,
    right=2.5cm,
    top=2.5cm,
    bottom=2.5cm
}
\usepackage{graphicx}
\usepackage{float}
\usepackage{hyperref}
\usepackage[english]{babel}
\usepackage{amsmath}
\usepackage{amssymb}
\usepackage{array}
\usepackage{booktabs}
\usepackage{caption}
\usepackage{subcaption}
\usepackage{enumitem}
\usepackage{fancyhdr}
\usepackage{titlesec}
\usepackage{url}
\usepackage{textcomp}
\usepackage{xcolor}
\usepackage{listings}
\usepackage{multirow}
\usepackage{longtable}
\usepackage{tabularx}
\usepackage{times}

% Set up page style
\pagestyle{fancy}
\fancyhf{}
\rfoot{\thepage}
\renewcommand{\headrulewidth}{0pt}

% Title formatting
\titleformat{\chapter}[display]{\normalfont\bfseries}{\chaptertitle\thechapter}{0pt}{\Huge}
\titlespacing*{\chapter}{0pt}{50pt}{40pt}

% Document info
\title{Template Document}
\author{Academic Support Agent}
\date{\today}

\begin{document}
\maketitle

\checkyears

\vspace*{5mm}

%\setkeys{Gin}{draft} % images non affich es
\setkeys{Gin}{draft=false} % images affich es

%\linenumbers % Numerotation des lignes pour relecture
%Select here the language 
\selectlanguage{french}   
%\selectlanguage{english} 


\def\author{Albert Tarantola}
\def\title{Le Rapport de Stage de Master}
\def\shorttitle{Rapport de stage}
\def\unit{IPGP}
\def\team{G omagn tisme \& Sismologie}
\def\spe{G ophysique ou autre} % Parcours
\def\supervisor{A. Fournier \& G. Moguilny}
\def\mydate{\today}
\def\keywords{Sciences de la Terre,  master, rapport, conseils}

\Entete


\begin{abstract}
Le stage de recherche du master est essentiel dans la formation d'un scientifique, 
car il permet un premier contact avec le monde de la recherche. 
La r daction du rapport est une composante importante du stage.
Il doit  tre r dig  en suivant le style litt raire habituel des articles 
scientifiques, avec ses usages bien codifi s. 
Les diff rentes parties du rapport : 
\emph{R sum , Introduction, R sultats, Discussion, Conclusion, Remerciements, 
R f rences bibliographiques,\/} et \emph{Annexes,\/} 
ont chacune des objectifs bien d \-ter\-mi\-n s. 
On doit chercher la perfection dans la forme.
\end{abstract}

\newpage

\newpage

\section{Introduction}

Le stage de recherche du master permet aux  tudiants d'apprendre le m tier 
de chercheur, en participant au travail de recherche d'un laboratoire. 
Le travail d'un chercheur ressemble beaucoup   celui d'un journaliste : 
il faut d'abord avoir l'id e d'un sujet d'enqu te, ensuite, il faut que 
l'enqu te fournisse des informations et, enfin, il faut la publier. 
Une grande quantit  d'enqu tes journalistiques, et de recherches 
scientifiques, sont rest es ignor es parce que leur publication a  t  
b cl e.

R diger un rapport scientifique prend beaucoup de temps. Apr s un 
``premier jet'', il faut souvent tout r crire. 
Lorsque la forme du rapport semble acceptable, il faut demander l'avis 
de ses coll gues, y apporter des modifications (souvent dramatiques), 
et it rer le processus
jusqu'  ce qu'un document de lecture facile ait  t  produit. Il~existe 
un \emph{style litt raire\/} tr s strict que les jeunes chercheurs ont parfois 
du mal   accepter. Si l'originalit  est indispensable pour le sujet de 
recherche, elle est simplement distrayante dans la publication,
o  on ne demande que de la perfection dans la pr \-sen\-ta\-tion.


\section{Les diff rentes parties d'un article}


\subsection{Le r sum }

Partie tr s importante du rapport.
Il doit permettre de se faire une id e pr cise des r sultats.
Le style est identique   celui de l'article, simplement en r sum .
Oubliez que le r sum  est publi  au d but de l'article et r digez 
comme s'il devait  tre la seule chose qui existe : ne faites pas 
r f rence   l'article (si vous disiez ``nous montrons que\dots'' 
vous seriez en train de le faire). 
Ne parlez pas de ``certaines'' hypoth ses ou r sultats ; dites lesquels 
ou lesquelles (ou n'en parlez pas si ce n'est pas important). Il est difficile de v hiculer beaucoup 
d'information en peu de lignes, mais c'est indispensable.


\subsection{L'introduction}

Sert   situer le contexte du travail. Pourquoi est-ce int ressant ? 
Qui a travaill  avant sur ce sujet? Quelle d marche avons-nous suivi ?
Il faut aussi dans l'introduction expliquer l'organisation de l'article 
(apr s quelques rappels de la th orie de base,  l'instrument 
utilis  est d crit dans la section 3.2\dots), sans que cela 
devienne une ``table des mati res''.


\subsection{La m thodologie, l'analyse, la discussion}

 videmment, la partie la plus importante dans le fond, tr s difficile 
   crire. C'est ici que la d marche scientifique doit  tre 
scrupuleusement suivie. De la rigueur, de la rigueur, et de la clart .


\subsection{Les conclusions}

Les conclusions doivent d j  appara tre dans le r sum  et le 
corps de l'article. Ici, en utilisant un langage technique ---permis par la
lecture de l'ensemble du rapport,--- on doit oublier les d tails, faire ressortir l'essentiel et \emph{ouvrir des perspectives\/}.


\subsection{Les r f rences bibliographiques}

Il existe plusieurs styles pour les r f rences bibliographiques. 
Il en existe m me trop : raison de plus pour ne pas en inventer davantage. 
Pour le rapport de stage du master, suivez le style sugg r  plus loin 
dans cet article. 

Notez que l'article de \cite{cao1990simultaneous} contient, dans 
la liste de r f rences bi\-blio\-gra\-phiques, le nom explicite de tous 
les auteurs.
L'article de \cite{crase1992nonlinear} est sorti.  
 Il y a aussi un article qui n'a  t  que 
soumis \citep{mosegaardsoumis}. 
 En fait, si 
vous utilisez \LaTeX, le formattage des r f rences peut  tre fait
automatiquement et vous n'avez pas   vous occuper du style 
(quand vous aurez besoin d'aide, adressez-vous   Genevi ve Moguilny,
au troisi me  tage). Pour plus de pr cisions sur 
\textsc{Bib}\TeX{} et \texttt{natbib}, 
 consultez  galement la partie~\ref{sec:natbib}. 

Lors de votre pr sentation orale, souvenez-vous que la locution latine 
``et al.'' est r serv e   l' crit : on ne dit pas ``Cao et al.'', 
mais ``Cao et collaborateurs''.


\subsection{Les annexes}

Un article doit  tre tr s lisible. Pour cela, il est utile de mettre en 
annexe toute d monstration ou discussion un peu lourde qui peut 
 tre distrayante. Il est pr f rable de n' crire dans l'article, 
dans un langage simple, que la trame des conclusions, qui seront justifi es 
dans le d tail seulement dans les annexes.


\section{Quelques remarques d'ordre g n ral}

L'information essentielle du rapport doit se trouver   {\bfseries quatre} endroits
diff rents : dans le r sum , dans le corps de l'article, dans les 
l gendes des figures, et dans la conclusion.

Il faut lire beaucoup d'articles scientifiques. C'est la seule fa on 
de saisir le \emph{style litt raire\/} d'un rapport scientifique.
N'envisagez pas d' crire votre rapport sans une bonne connaissance des 
revues scientifiques importantes (\emph{Nature, Science, Journal of 
Geo\-phy\-si\-cal Research, Geophysical Journal International, Geophysics, 
Physics of the Earth and Planetary Interiors, Earth and Planetary Science 
Letters\dots\/}). Le style d'un rapport scientifique n'est pas celui des 
revues de vulgarisation (\emph{Science et Vie, La Recherche\dots\/})
qui, par ailleurs, peuvent  tre excellentes.

De plus, pour ce rapport, au format A4 :
\begin{itemize}
\item la premi re page doit contenir le r sum  de 20   25 lignes, 
\item la deuxi me page doit contenir le Sommaire,
\item le nombre de pages ne doit pas exc der 30,
\item les pages doivent  tre  num rot es en haut   droite/gauche,
\item le nombre de pages doit  tre pair (si les rapports devaient  tre concat n s  pour impression).
\end{itemize}

\subsection{Les figures}

Si possible, les points clefs de l'article doivent  tre illustr s par une 
figure (celle qui vaut plus que 100~mots).

En utilisant des logiciels habituels de dessin (e.g., {\sl Illustrator\/}, 
 {\sl gimp\/}), 
il est possible de produire des figures propres qui peuvent s'int grer 
dans l'article en utilisant un des \textit{packages\/} existants 
comme \texttt{graphicx}.
Essayez, dans la mesure du possible, de produire des figures en format 
\texttt{pdf} (\textit{Portable Document Format\/}).

Il n'est pas interdit d'utiliser une figure extraite d'un livre ou d'un 
article (en la passant au scanner), mais il ne faut pas oublier d'en donner 
la r f rence exacte.

%Pendant que le rapport est en cours d' laboration, il peut  tre commode 
%de situer toutes les figures   la fin du texte (\begin{verbatim}
%\usepackage{endfloat} \end{verbatim}). 
%Mais 
Dans la forme finale du rapport, les figures doivent  tre   leur place.
Les l gendes des figures doivent  tre assez explicites pour permettre 
de comprendre le travail, m me si on ne lit que le r sum  et les 
l gendes des figures.

L'inclusion d'une image simple est montr e sur la figure \ref{logo},
et un positionement plus sophistiqu  sur les figures \ref{gmt}   \ref{draftimage}.

\begin{figure}
\centerline{\includegraphics[width=0.4\linewidth]{LogoIPGPUPC}}
\caption{Logo de l'IPGP depuis juin 2022.}
\label{logo}
\end{figure}


\begin{minipage}{\linewidth}
Deux erreurs sont tr s communes 
\begin{enumerate}
\item Des figures inutilement grandes : la r solution d'une figure 
est plus importante que la taille et souvent une m me figure est 
plus ``sympathique'' en plus petit.
\item Des figures tourn es de 90$^\circ$ qui forcent le lecteur   
tourner la page, pour l'utiliser en mode ``paysage''. 
C'est tr s rare qu'il existe des bonnes raisons   cela.
\end{enumerate}
\end{minipage}

\begin{figure}

\begin{minipage}[b]{0.48\linewidth}
\includegraphics[width=\linewidth]{gmt1}
\caption{Image produite par \href{http://gmt.soest.hawaii.edu/}{GMT}.}
\label{gmt}
\end{minipage}
\hfill
\begin{minipage}[b]{0.48\linewidth}
\fboxsep 0pt
\centerline{\fbox{\includegraphics[bb=177 146 200 161,clip=,width=0.75\linewidth]{gmt1}}}
\caption{Zoom sur la Fig~\ref{gmt}.}
\label{zoomimage}
\end{minipage}

\vspace*{5mm}

\begin{minipage}{0.48\linewidth}
\includegraphics[width=\linewidth,draft]{gmt1}
\end{minipage}
\hfill
\begin{minipage}[b]{0.48\linewidth}
\caption{Comme la Fig~\ref{gmt}, mais en \texttt{draft}.}
\label{draftimage}
\end{minipage}

\end{figure}


\subsection{Les  quations}

 crivons quelques lignes de texte technique pour illustrer la fa on 
normale d' crire des  quations.

Soit un milieu parfaitement  lastique, et soient
\ $ \sigma^{ij} ({\bf x},t) $ \ 
et
\ $ \varepsilon^{ij} ({\bf x},t) $ \ 
respectivement les composantes 
\ $ \{i,j\} $ \ 
de la contrainte et de la d formation au point 
\ $ {\bf x} $ \ 
et   l'instant 
\ $ t $~.
Pour des petites d formations, nous pouvons  crire, au premier ordre,
\begin{equation}
     \sigma^{ij} ({\bf x},t) 
     \ = \  
     \sigma_0^{ij} ({\bf x},t) 
     + 
     c_0^{ijkl} ({\bf x}) 
     \,  
     \varepsilon^{kl} ({\bf x},t) 
     \ , 
          \label{eq:loiHooke}
\end{equation}
o  le terme 
\ $ \sigma_0^{ij} ({\bf x},t) $ \ 
est appel  \emph{pr contrainte\/}, et o  les coefficients 
\ $ c_0^{ijkl} ({\bf x}) $ \ 
sont les \emph{constantes  lastiques\/}. 
L' quation~(\ref{eq:loiHooke}) repr sente la \emph{ Loi de Hooke\/}.
Maintenant, que repr sente l' quation
\begin{equation}
     E({\bf x},t) 
     \ = \  
     \frac{1}{2} 
     \, 
     \sigma^{ij} ({\bf x},t) 
     \, 
     \varepsilon^{ij} ({\bf x},t)
     \ ? 
          \label{eq:energydensity}
\end{equation}
Peu importe en v rit , car elle est hors sujet. 
Par contre, en utilisant la loi de Hooke~(\ref{eq:loiHooke}) et le principe
 d' quivalence entre la masse inerte et la masse pesante, 
on d montre sans peine l'\emph{ quation des ondes  lastiques\/} :
\begin{equation}
     \rho \ 
     \frac{\partial^2 u^i}{\partial t^2} 
     - 
     \frac{\partial}{\partial x^j} 
     \left( 
           c_0^{ijkl} 
           \, 
           \frac{\partial u^k}{\partial x^l} 
     \right)
     \ = \  
     0 
     \ . 
          \label{eq:ondes}
\end{equation}

La discussion ci-dessus illustre quelques unes des r gles   utiliser dans 
un texte contenant des  quations :

\begin{itemize}
     \item  Il faut introduire dans le texte {\bfseries toutes} les variables.

     \item  Il faut num roter {\bfseries toutes} les  quations.

     \item  Dans la mesure du possible, il faut donne un \emph{nom\/}   
            chaque  quation, et, dans la suite du texte, dire 
            ``l' quation des ondes~(\ref{eq:ondes})'' 
            au lieu de ``l' quation~(\ref{eq:ondes})''.

     \item  Il faut \emph{ponctuer\/} les  quations, car elles font 
            partie du texte. Ici, non seulement nous avons des  quations 
            se terminant par des virgules et des points, mais m me 
            une  quation se terminant par un point d'interrogation.

     \item  N'abusez pas du ``~:~'' pour arriver aux  quations : 
            dans les exemples ci-dessus, seulement une  quation le permet.
\end{itemize}

\subsection{La typographie}

Soignez l'orthographe.  vitez m me les erreurs comme celle de ne pas accentuer 
les majuscules (cette phrase commence par une majuscule accentu e, et pour vous
 con\-vain\-cre, regardez les couvertures de livres ou les titres des
journaux, ou v rifiez dans votre dictionnaire que l'on  crit bien 
 tats-Unis\dots).

\subsection{\LaTeX, Word, ou quoi d'autre?}

Tout logiciel de pr paration de documents peut  tre utilis  pour produire 
votre rapport. \LaTeX \ est bien adapt    la publication scientifique, 
avec ses notations ma\-th \-ma\-ti\-ques. Mais, surtout, il est devenu le 
standard pour la publication technique.

Pour votre rapport vous devrez utiliser exactement le style illustr  par ce document. 
Si vous utilisez \LaTeX, vous n'aurez rien de sp cial   faire 
(ni   d finir les fontes, ni les espaces entre sections, etc.), car 
ceci n'est qu'une adaptation mineure du style ``article'' de \LaTeX.
Je d conseille de choisir un autre logiciel, moins standard ou moins technique.
Si vous \^{e}tes oblig s de le faire, reproduisez exactement le style de ce 
document. Pensez   bien utiliser des caract res en ``11~points'' 
\footnote{Au passage, pour  viter de couper entre un nombre et son unit , 
utiliser le tilde (blanc ins cable) : \texttt{11$\sim$points}.} 
pour le texte principal (la taille de ces caract res). 

\subsection{Les erreurs courantes     viter}

Une erreur courante des auteurs inexp riment s est de ne pas donner assez
 d'importance aux r f rences bibliographiques : r f rences manquantes, 
r f rences cit es dans le texte, mais absentes de la liste finale, etc. 
Souvent aussi, la liste des r f rences ne suit pas un des formats habituels, 
ou il y manque un des  l ments (ann e, page, etc.).

Souvent les figures sont trop grandes, ou bien elles d passent des marges de 
l'article. 
 vitez d'obliger le lecteur   tourner la revue : il est rare qu'une figure
 doive absolument avoir un format paysage sur une page enti re.
Parfois il y a une multiplicit  de l gendes : on  crit des phrases dans la
 figure elle-m me.      

Il arrive que le r sum  ou l'introduction manquent, ou que l'on confonde 
leurs r les respectifs.     

Si un document a un seul auteur, on ne doit jamais utiliser le ``nous'',  
qui est r serv  au Pape. Le ``je'' doit  tre utilis  dans des tr s 
rares occasions. Soyez neutre : ``il'' est possible de montrer que\dots 

Une grave erreur consiste   sous-estimer le temps qu'il faut pour
 produire un rapport sous sa forme finale : commencez   r diger d s la mi-mai.
M me si vos recherches sont loin d'avoir abouti, vous pouvez d j  commencer
   mettre la bibliographie ---\,et, surtout, vos id es\,--- en place.
 
 
 Autres erreurs communes :
\begin{itemize}
\item utiliser le retour chariot (\verb"\\") pour finir un paragraphe,
\item essayer de forcer le placement des objets flottants comme les figures et les tables (option \verb"[h]"),
\item continuer   saisir du texte, sans corriger les erreurs de compilations, et ne pas regarder les causes d'erreurs dans les fichiers log.
\end{itemize}



\subsection{Utilisation de \texttt{natbib} et \texttt{bibtex} pour les r f rences}
\label{sec:natbib}
Il est imp ratif que vous preniez l'habitude d'utiliser \textsc{Bib}\TeX{} (ou \textsc{Bib}\LaTeX{})
en conjonction avec le \textit{package\/} \texttt{natbib} pour g rer votre
 base de donn es de r f rences et son adressage dans votre texte. L'utilisation de \textsc{Bib}\TeX{} 
est 
 d taill e dans le texte de \cite{moguilny2021introduction}. Ne pas oublier d'ajouter les DOI 
 (\textit{Digital Object Identifier}) dans les entr es de la base de donn es \texttt{.bib}.

 Quelques emplois possibles :

 \verb"\cite{fournier2010introduction}"  donnera \cite{fournier2010introduction}  \\
 \verb"\citep{fournier2010introduction}" donnera \citep{fournier2010introduction}  \\
 \verb"\citep[voir par exemple][\S 2]{fournier2010introduction}"  donnera \\
 \hspace*{1cm}\citep[voir par exemple][\S 2]{fournier2010introduction}  \\
 \verb"\citeauthor{fournier2010introduction}"  donnera \citeauthor{fournier2010introduction}  \\
 \verb"\citeyear{fournier2010introduction}"  donnera \citeyear{fournier2010introduction}  

 La documentation compl te de \texttt{natbib} en fran ais est disponible   l'adresse \\
\href{http://mirrors.ircam.fr/pub/CTAN/info/translations/natbib/fr/f-natbib.pdf}
{http://mirrors.ircam.fr/pub/CTAN/info/translations/natbib/fr/f-natbib.pdf}

\section{Conclusion}

 crire un rapport de stage de master n'est pas ais . 
Cela prend du temps, et les difficult s sont souvent sous-estim es.
Mais cette r daction est une composante essentielle du stage de recherche.

Le \emph{style\/} d'un rapport technique est bien codifi , et laisse peu de place
  l'originalit , qui, par contre, aura libre cours dans le choix du sujet
ou des m thodes de recherche. On demande une grande \emph{ clart \/} et une grande 
\emph{concision\/} dans la pr sentation. 
Pour cela, il faut s'y prendre avec du temps : un document  crit se travaille 
beaucoup avant de le divulguer.

\section{Remerciements}

G n ralement, vous serez le seul signataire du rapport ---\,et, donc, le seul 
responsable\,---.
Votre directeur de stage n'appara tra que dans cette section 
``\emph{Remerciements\/}''. 
Si votre rapport contient des r sultats scientifiques, veillez   faire justice, 
et faire appara tre ici les apports respectifs de votre directeur de stage
et de vous m me.

Un style typique pour cette section est le suivant :

\paragraph*{Remerciements.} Ce stage de master a  t  effectu  dans le 
Laboratoire de G o\-pan\-tou\-flis\-me de 
l'Institut de Physique du Globe de Paris, sous la direction de Napol on Bonaparte, 
que je voudrais remercier pour la qualit  du caf  qu'il sait faire. 
 mile Zola m'a aid  avec le logiciel \LaTeX \ de pr paration de documents. 
Le logiciel de dessin utilis  appartient   l'Universit  de Cochabamba.
Cette recherche a  t  men e   bien gr ce   une bourse du 
Minist re des Anciens Ministres, et le microscope aquatique utilis  a 
 t  financ  par le Centre National de la Recherche Parapsychologique. 
 
\newpage

\addcontentsline{toc}{section}{R f rences}

%\nocite{*} % pour afficher toutes les entr (c)es (m me celles non cit es)
            % de la base de donn es bibliographique
            
\newpage


\section{Annexes}
 
\subsection{Pr cisions sur un certain point}
  
\subsection{Pr cisions sur un autre point}
\end{document}